\part{Context modelling and application traffic generation}

\section{Context engine}

\begin{qa}{165}{What does context mean in CARS-SF's mathematical and operational model?}
A discrete variable $c_t\in\mathcal{C}$ = \{normal, school, harvest, disaster, medical\_emergency, fiber\_cut,
backhaul\_degradation, core\_switch\_failure\}, inferred each tick by
\[
c_t=f(\text{clock}_t,\ \text{calendar},\ A_t,\ o_t),
\]
where $A_t$ is the set of active alerts and $o_t$ the admin override.
Operationally, $c_t$ selects the priority weights $w(c_t)$ and the design demand $d(c_t)$ from
\code{scenarios.yaml}.
\end{qa}

\begin{qa}{166}{What context variables are represented?}
\begin{itemize}
\item time of day and weekday;
\item month (harvest months 3, 4, 10, 11);
\item school days and hours (Mon--Fri, 08--15);
\item harvest daylight hours (06--18);
\item alerts (flood, fire, earthquake, disaster, medical\_emergency, fiber\_cut);
\item an admin override.
\end{itemize}
Source: \code{registries/timelines.yaml}, \code{context/engine.py}.
\end{qa}

\begin{qa}{167}{Which are static, slowly varying, event-driven or continuously changing?}
\begin{itemize}
\item Static: the calendar configuration.
\item Slowly varying: season (month).
\item Continuously changing: time of day.
\item Event-driven: alerts and the admin override.
\end{itemize}
Network condition is \emph{not} context; it is handled by detection and healing.
\end{qa}

\begin{qa}{168}{How are emergencies, congestion, link failures, holidays and power failures represented?}
\begin{itemize}
\item Emergencies: alerts mapped to contexts (flood, fire and earthquake $\to$ disaster).
\item Link failures: scenarios with a \code{failure} spec driving the injector.
\item Congestion: emerges from offered load; it is not a context.
\end{itemize}
\honest{Holidays and power failures exist only in the RSNSD generator's 19 scenarios, not in the live CARS-SF context
registry.}
\end{qa}

\begin{qa}{169}{Why is a scenario registry better than hard-coded conditions spread through the code?}
\why{One definition of priorities and demand per context for the allocator, the traffic, the RL environment and the
analysis.}
\choice{\code{scenarios.yaml}, read everywhere.}
\alt{\code{if context == ...} branches (the original \code{DemandGenerator} did this, and it diverged:
\code{medical_emergency} silently fell back to the defaults).}
\trade{Consistency and a new context added by editing data; the cost is no schema validation of the YAML (Q126).}
\ev{The bug that motivated this was found in the verification audit.}{\Lthree}
\end{qa}

\begin{qa}{170}{What are the context engine's inputs and outputs?}
In: tick, timeline events (alerts, clear, admin), clock start and minutes per tick.
Out: \code{\{tick, clock, context, reason, alerts, changed\}} per tick, plus history.
\end{qa}

\begin{qa}{171}{Does context describe the environment, the application's urgency, the network's condition, or all three?}
The environment and the \emph{societal} urgency it implies (school is in session, a flood is happening). Application
urgency is in the per-app SLAs; network condition is in telemetry. Keeping them separate prevents a congested link
from being read as ``an emergency''.
\end{qa}

\begin{qa}{172}{How do you distinguish a change in demand from a change in capacity?}
Demand is measured where traffic is still unshaped: access-port rx counters, and HTB sent+dropped. Capacity changes
are healing states (backup at 50\,Mb/s, core down) that enter as constraints. The predictor adds a discriminator:
rising RTT \emph{with} local queue growth or drops is congestion (demand); rising RTT \emph{without} it is the
backhaul (capacity or quality).
\end{qa}

\begin{qa}{173}{How does the system prevent two context sources from issuing conflicting priorities?}
A strict precedence produces exactly one context: admin override $>$ alerts (medical $>$ disaster $>$ fiber cut) $>$
fixed scenario $>$ school hours $>$ harvest $>$ normal (\code{ContextInferenceEngine.infer}; unit-tested).
\end{qa}

\begin{qa}{174}{How are simultaneous events combined?}
They are not blended. Precedence picks one context.
\honest{Flood plus school hours yields disaster weights, with no ``school during a flood'' mix. Blending (e.g.\ a max or
weighted combination of weight vectors) is possible future work.}
\end{qa}

\begin{qa}{175}{What happens if an emergency begins while the network is already congested?}
The context switch triggers an immediate re-plan with emergency weights and the emergency's design demand as the
prior, because the estimator restarts with warm-up. The floors still hold. Congestion boosts from the healer multiply
on top. In the timeline runs, disaster begins mid-run under load.
\end{qa}

\begin{qa}{176}{What happens if an emergency ends while the system is still recovering from a link failure?}
Context and healing are independent state machines. The controller re-plans for the new context, and the healing
constraints (backup ceiling) still apply until fail-back.
\honest{This exact overlap is not in the test matrix: the timelines do not inject faults.}
\end{qa}

\begin{qa}{177}{How do you model event duration and transitions?}
Timeline events carry a raise tick and a clear tick (flood from tick 12 to tick 24). The calendar produces
deterministic transitions. The RL environment uses Markov switching ($p_{stay}=0.92$).
\end{qa}

\begin{qa}{178}{How do you avoid abrupt policy oscillations when context changes rapidly?}
\begin{itemize}
\item Contexts change only on events or calendar boundaries.
\item The controller applies a change only if a slice moves by at least 1\,Mb/s.
\item Healing boosts have a cooldown.
\end{itemize}
\honest{There is no dwell-time hysteresis on context itself. A flapping alert feed would flap priorities; a minimum
context dwell time is a one-line extension.}
\end{qa}

\begin{qa}{179}{Is context categorical, numerical, probabilistic, or a combination?}
Categorical: one id, one-hot in the RL observation. The calendar inputs are numerical, but inference is rule-based and
deterministic, with no probabilities.
\end{qa}

\begin{qa}{180}{How are unknown or unsupported context values handled?}
An unknown \emph{alert} maps to disaster (the conservative default). An unknown context id raises in
\code{get_context}, and the orchestrator's demand lookup falls back to normal. An unknown timeline raises before the
run.
\end{qa}

\begin{qa}{181}{How would noisy or incorrect context affect the allocation decision?}
Wrong weights shift shares between slices, but demand-capping limits the waste: a slice given a high weight but low
demand gets only its demand, and the surplus is redistributed. The floors bound starvation. The \code{no_context}
ablation (always ``normal'' weights) changed school utility from \NAblSchoolFullUtil{} to \NAblSchoolNoContextUtil{}, i.e.\ no
measurable damage. In this setting even fully wrong context is harmless; equally, correct context is not what drives
outcomes.
\end{qa}

\begin{qa}{182}{How do you validate that a simulated context produces the intended downstream behaviour?}
The integration test \code{test_timeline_switches_context_and_traffic} asserts that school and disaster occur and that
the controller re-planned with reason ``context''. The timeline study's \code{feedback.csv} shows per-tick context,
allocation and throughput (Fig.~\code{timeline_throughput}).
\end{qa}

\begin{qa}{183}{How would real context enter the system?}
\begin{itemize}
\item Operator UI (admin override).
\item Alert feeds such as the Common Alerting Protocol from civil protection.
\item School timetables and agricultural calendars from local institutions.
\item Application metadata (a telemedicine session start) as an alert-like event.
\end{itemize}
The engine already exposes \code{raise_alert}, \code{clear_alert} and \code{set_admin_override}.
\end{qa}

\begin{qa}{184}{What are the security and trust implications of externally supplied context?}
A spoofed ``disaster'' alert could deprioritise everything except healthcare. Mitigations:
\begin{itemize}
\item authenticated or signed feeds;
\item allow-listed sources;
\item rate limits and minimum dwell time;
\item audit logs (the history is kept);
\item a bounded blast radius, because the floors still hold.
\end{itemize}
\honest{None of the authentication is implemented.}
\end{qa}

\section{Traffic and workload modelling}

\begin{qa}{185}{Why generate application-specific traffic instead of one generic source?}
\why{Slicing and classification are only meaningful if services differ in burstiness, direction, rate and transport
the way real applications do.}
\choice{Seven models (VBR video, VoIP, IoT periodic, Pareto on/off, paced bulk TCP, web request/response, DASH) across
16 applications (\code{traffic/models.py}, \code{traffic/generator.py}).}
\alt{iperf3 CBR (the original code), or trace replay.}
\trade{Realistic mixes and ground-truth labels; the cost is that the parameters are assumed (Q188).}
\ev{Unit tests check rate conformance and I-frame bursts.}{\Ltwo}
\end{qa}

\begin{qa}{186}{Which generators are used per domain?}
\begin{itemize}
\item \textbf{Healthcare}: telemedicine VBR video (up, 4\,Mb/s), ECG periodic (50\,pkt/s, 200\,B), emergency voice
(VoIP talkspurts), ambulance on/off (1.5\,Mb/s).
\item \textbf{Education}: live classroom VBR video (down, 20), library bulk TCP (down, 40), exam web (3), LMS web (8).
\item \textbf{Agriculture}: soil (every 2\,s), weather (every 5\,s), drone video (up, 8), irrigation on/off (0.2).
\item \textbf{General}: web (4), DASH streaming (15), social (3), banking (0.5).
\end{itemize}
All rates are nominal and are scaled to each context's demand.
\end{qa}

\begin{qa}{187}{How are packet sizes, inter-arrivals, durations, protocols and offered loads determined?}
Per-model parameters in \code{applications.yaml}:
\begin{itemize}
\item video: fps, GOP, I/P ratio 4, lognormal coefficient of variation 0.25;
\item VoIP: 172\,B at 50\,pkt/s, talkspurt exponential means 1.0/1.35\,s;
\item IoT: period and size with $\pm$jitter;
\item on/off: Pareto shape 1.5;
\item web: Pareto response sizes, shape 1.3.
\end{itemize}
Offered load: each service's nominal rates are scaled so they sum to the context demand, then split across UEs.
Durations last for the run.
\end{qa}

\begin{qa}{188}{Which parameters are from published evidence, which are measured, and which are assumptions?}
\begin{itemize}
\item Standard-derived: G.711-style voice framing (50\,pkt/s, 160\,B payload); typical video frame rates and GOP structure.
\item Classic models from the literature: Pareto on/off for self-similar traffic\textsuperscript{†}; heavy-tailed web sizes.
\item Measured: nothing from real rural networks.
\item Assumptions: all rates and the per-context demands.
\end{itemize}
\end{qa}

\begin{qa}{189}{How do you model bursty traffic?}
Pareto on/off (heavy-tailed bursts at an elevated rate so the mean is preserved), VBR video I-frames sent as 1\,ms
bursts, and web responses with Pareto sizes.
\end{qa}

\begin{qa}{190}{How do you model periodic telemetry?}
\code{iot_periodic}: fixed size every \code{period}$\times(1\pm\text{jitter})$. Demand scaling adds
\emph{devices} (instances), not rate.
\end{qa}

\begin{qa}{191}{How do you model sustained video or classroom traffic?}
GOP-structured VBR. The I-frame is 4$\times$ a P-frame, frame sizes are lognormal, and each frame is split into
$\le$1200\,B packets in a 1\,ms burst at 25--30\,fps. DASH streaming is TCP chunk downloads every 2\,s.
\end{qa}

\begin{qa}{192}{How do you model emergency voice or medical streams?}
Voice: 50\,pkt/s during exponential talkspurts. ECG: periodic 200\,B every 20\,ms. Telemedicine: VBR video uplink.
\end{qa}

\begin{qa}{193}{How do you represent a sudden increase in demand?}
A context switch re-spawns generators with the new demand scaling (timelines), and the load sweep (0.5--1.5$\times$).
Within a context, bursts come from the models themselves.
\end{qa}

\begin{qa}{194}{How do you model correlated bursts across several services?}
Only through context switches, where all services change together.
\honest{There is no explicit cross-service burst correlation (a copula or shared modulating process). Independent
seeds make bursts independent, which probably flatters statistical multiplexing.}
\end{qa}

\begin{qa}{195}{How do you ensure generated workloads are reproducible?}
Every generator gets a deterministic seed (\code{seed·100003 + app·101 + ue}), the schedules are pure functions of
it, and a test confirms identical schedules for identical seeds. Realised packet timing still varies with OS
scheduling.
\end{qa}

\begin{qa}{196}{How are random seeds managed across multiple generators?}
The run seed comes from the matrix. Per-application, per-UE seeds are derived from it; instances add $7919\cdot i$.
A context switch increments the run seed so new traffic is not a replay.
\end{qa}

\begin{qa}{197}{How do you validate that a generator produces the distribution it claims?}
\code{tests/test_traffic_models.py} checks the long-run rate against the nominal over 120\,s (CBR 2\%, video 8\%,
IoT 5\%, on/off and VoIP 35\% because of heavy tails), I-frame bursts above 3$\times$ the median frame, ordering and
bounds. At run time, the dataset shows delivered rates per application.
\end{qa}

\begin{qa}{198}{How do you avoid unrealistic traffic that makes classification artificially easy?}
\begin{itemize}
\item Ports, DSCP, TTL and timestamps are excluded from features.
\item Evaluation holds out runs and UEs.
\item Traffic is multi-model within each service.
\end{itemize}
\honest{The four services' model mixes are still ours. The F1 of \NSnRfFone{} measures separability of
\emph{our} generators, not of real services. Real-trace validation is needed.}
\end{qa}

\begin{qa}{199}{How do you distinguish offered load from delivered throughput?}
Offered: the generator schedule, the access-port rx and HTB sent+dropped. Delivered: bytes received at the
destination endpoint (matched packets). The gap is loss.
\end{qa}

\begin{qa}{200}{How do you tell whether the bottleneck is the sender, receiver, virtual link, queue or path?}
\begin{itemize}
\item Delay floors equal configured propagation exactly (\NFloorEdge/\NFloorCore\,ms), so excess delay is queueing.
\item HTB class backlog and drops locate the queue.
\item Generator CPU (\NGenCpuAfter\% per 40\,Mb/s) rules out the sender.
\item In the 5G case, comparing UPF-shaped and unshaped modes exposed UERANSIM's user-space GTP as a bottleneck
(Q598f).
\end{itemize}
\end{qa}

\begin{qa}{201}{Why not replay real packet traces?}
No rural service traces with service labels exist that we can access. Replay also cannot react to slicing (TCP
replays do not back off), and labels and privacy are hard. Synthetic, closed-loop generators react to the network.
\end{qa}

\begin{qa}{202}{Benefits and drawbacks of synthetic workloads vs real traces?}
Benefits: control, ground-truth labels, reproducibility, scaling to any context, closed-loop TCP behaviour.
Drawbacks: realism is unproven, unknown unknowns are missing, and classification may be optimistic.
\end{qa}

\begin{qa}{203}{How do you avoid training and evaluating on effectively identical generated flows?}
By group hold-out: test windows come from runs and UEs never seen in training (\NSnGroups{} groups). Seeds differ
across runs.
\honest{The same \emph{generator family} appears in train and test. Generalisation to new applications is untested.}
\end{qa}

\begin{qa}{204}{How would you test generalisation to a previously unseen application?}
Leave-one-application-out: train without, e.g., \code{general.social_media}, test on it, and measure whether it
maps to the right service. Expect degradation. Add an ``unknown'' rejection threshold (Q396).
\end{qa}

\begin{qa}{205}{What happens if a traffic generator fails midway?}
Its log (\code{generated/logs/<run>/app@host.log}) records the error, and the flow disappears from later windows.
The validator's class-coverage score drops if a whole service vanishes. Nothing restarts it.
\honest{An integration assertion on per-app record counts would catch this; it is only implicit today.}
\end{qa}
